\begin{ouverture}[college1]{2}{L’air, les gaz et la pression}{Partie 1 · Organisation et transformations de la matière}{Pourquoi une seringue bouchée résiste-t-elle ?}
\objectif{Décrire l’air comme un mélange de gaz.}
\objectif{Interpréter compression et détente avec un modèle particulaire.}
\objectif{Distinguer pression, masse et volume.}
\prerequis{Corps pur, mélange et volume.}
\end{ouverture}
\pcPage{Le cours}{L’air, les gaz et la pression}
\phantomsection\label{ess-2}
\essSection{1}{De quoi l’air est-il fait ?}
L’air sec contient environ $78\,\%$ de diazote, $21\,\%$ de dioxygène et $1\,\%$ d’autres gaz en volume. L’air réel contient aussi une quantité variable de vapeur d’eau. Ces constituants sont mélangés ; ils ne forment pas des couches séparées dans une seringue. Le dioxygène permet de nombreuses combustions, mais l’air n’est pas du dioxygène pur.\par
\pcFigure{figures/c02-air.pdf}{Pour 100 parts d’air sec : proportions arrondies des principaux constituants.}
\essSection{2}{Un gaz occupe un volume}
Un gaz occupe tout l’espace disponible d’un récipient. Dans une seringue fermée, enfoncer doucement le piston réduit le volume sans retirer de gaz : c’est une compression. Dans le modèle, les molécules restent présentes et plus rapprochées ; elles ne deviennent pas plus petites. Lors d’une détente, l’espace disponible augmente.\par
\clearpage
\essSection{3}{La pression et ses unités}
Les particules d’un gaz exercent des actions sur les parois. Un capteur mesure la pression $p$, en pascals (Pa). On utilise aussi l’hectopascal : $1\unite{hPa}=100\unite{Pa}$, et le bar : $1\unite{bar}=100\,000\unite{Pa}$. À quantité de gaz et température constantes, une compression augmente la pression. On ne suppose pas cette température constante pour une compression brusque.\par
\pcFigure{figures/c02-pression.pdf}{Même quantité de gaz, volume plus faible : le modèle représente des particules plus rapprochées.}
\essSection{4}{L’air possède une masse}
Peser un ballon avant puis après retrait d’une partie de son air montre une différence de masse, si l’on maîtrise les conditions de mesure. À environ $20\,^{\circ}\mathrm C$ et à pression voisine de la pression atmosphérique, on retient comme ordre de grandeur $1{,}2\unite g$ par litre d’air. Cette valeur dépend fortement de la température et de la pression. Une indication en litres ne donne donc pas directement une masse universelle.\par
\begin{pcEncadre}{method}{Un exemple expliqué}Un échantillon de $500\unite{mL}$ d’air sec correspond approximativement à $0{,}21\times500=105\unite{mL}$ de dioxygène et $0{,}78\times500=390\unite{mL}$ de diazote, aux mêmes température et pression. Les $5\unite{mL}$ restants représentent les autres gaz. Les proportions restent inchangées lors d’un simple prélèvement homogène.\end{pcEncadre}
\pcPage{Expérimenter et retenir}{L’air, les gaz et la pression}
\phantomsection\label{activite-2}
\begin{pcEncadre}{method}{Investigation en classe}\textbf{Question.} Comment varie la pression quand le volume diminue ? Utiliser une seringue et un capteur scolaires, sans aiguille. Fermer le dispositif et attendre l’équilibre thermique. Relever la pression pour trois volumes autorisés par le professeur, sans forcer. Revenir doucement au volume initial pour tester la reproductibilité. \textbf{À discuter :} quantité de gaz, température, étanchéité et différence entre observation qualitative et relation mathématique. Ne pas essayer de provoquer l’éjection du bouchon.\end{pcEncadre}
\pcSousTitre{Ce que je dois savoir expliquer}\begin{questions}
\item L’air sec est un mélange dont le diazote et le dioxygène sont les constituants principaux.
\item Une compression diminue le volume disponible.
\item La masse d’un gaz enfermé se conserve si le système est étanche.
\item La pression se mesure avec un capteur et s’exprime en Pa.
\item Les propriétés d’un gaz dépendent des conditions de température et de pression.
\end{questions}
\begin{pcEncadre}{attention}{Une erreur à éviter}Une molécule ne « gonfle » pas pendant une détente : c’est principalement l’écartement entre les molécules qui change dans le modèle du gaz.\end{pcEncadre}
\pcRepere{cor-activite-2}{Pistes pour l’investigation}\pcGrille
\clearpage\markright{Je démarre}
\rubrique{Je démarre}
\par\noindent\begin{minipage}{\linewidth}
\begin{exercice}[1]{Les deux principaux gaz}
\begin{questions}
\item Nommer les deux principaux gaz de l’air sec.
\item Donner leurs proportions arrondies.
\end{questions}
\end{exercice}
\end{minipage}
\par\noindent\begin{minipage}{\linewidth}
\begin{exercice}[1]{Air ou dioxygène ?}
\begin{questions}
\item L’air est-il un corps pur ?
\item Le dioxygène est-il synonyme d’air ?
\end{questions}
\end{exercice}
\end{minipage}
\par\noindent\begin{minipage}{\linewidth}
\begin{exercice}[1]{Compresser}
\pcFigureEx{figures/ex-compression.pdf}{}
\begin{questions}
\item Le volume d’une seringue fermée passe de $60$ à $40\unite{mL}$. Comment nommer l’opération ?
\item Des molécules ont-elles été retirées ?
\end{questions}
\end{exercice}
\end{minipage}
\par\noindent\begin{minipage}{\linewidth}
\begin{exercice}[1]{Des unités de pression}
\begin{questions}
\item Convertir $1\,020\unite{hPa}$ en Pa.
\item Convertir $2\unite{bar}$ en Pa.
\end{questions}
\end{exercice}
\end{minipage}
\par\noindent\begin{minipage}{\linewidth}
\begin{exercice}[1]{Dans cent millilitres}
\begin{questions}
\item Quel volume correspond approximativement au dioxygène dans $100\unite{mL}$ d’air sec ?
\item Quel volume correspond aux autres gaz ?
\end{questions}
\end{exercice}
\end{minipage}
\clearpage\markright{Je m’entraîne}
\rubrique{Je m’entraîne}
\par\noindent\begin{minipage}{\linewidth}
\begin{exercice}[2]{Un prélèvement plus grand}
\begin{questions}
\item Dans $2\unite L$ d’air sec, calculer le volume correspondant au dioxygène.
\item Exprimer celui du diazote en litres.
\end{questions}
\end{exercice}
\end{minipage}
\par\noindent\begin{minipage}{\linewidth}
\begin{exercice}[2]{Lire un capteur}
\begin{questions}
\item Une mesure passe de $1{,}0$ à $1{,}4\unite{bar}$ pendant une compression lente. La pression augmente-t-elle ?
\item La masse augmente-t-elle forcément ?
\end{questions}
\end{exercice}
\end{minipage}
\par\noindent\begin{minipage}{\linewidth}
\begin{exercice}[2]{Une masse d’air}
\begin{questions}
\item On prend $1{,}2\unite{g/L}$ pour l’air dans les conditions du cours. Estimer la masse de $5\unite L$.
\item Nommer deux conditions à préciser.
\end{questions}
\end{exercice}
\end{minipage}
\par\noindent\begin{minipage}{\linewidth}
\begin{exercice}[2]{Peser par différence}
\begin{questions}
\item Un ballon équipé de sa valve passe de $425{,}6$ à $423{,}2\unite g$ après retrait d’air. Quelle masse a été retirée ?
\item Estimer son volume avec $1{,}2\unite{g/L}$.
\end{questions}
\end{exercice}
\end{minipage}
\par\noindent\begin{minipage}{\linewidth}
\begin{exercice}[2]{Modéliser une détente}
\pcFigureEx{figures/ex-detente.pdf}{À droite, compléter le modèle après détente sans fuite.}
\begin{questions}
\item On passe de 20 à 40 mL sans fuite. Comment modifier un dessin de particules ?
\item Que doit-on conserver dans le dessin ?
\end{questions}
\end{exercice}
\end{minipage}
\clearpage\markright{Je raisonne}
\rubrique{Je raisonne}
\par\noindent\begin{minipage}{\linewidth}
\begin{exercice}[3]{Une conclusion excessive}
\begin{questions}
\item Après une seule compression, un élève affirme que doubler la pression divise toujours le volume par deux. Cette mesure suffit-elle ?
\item Quel paramètre faudrait-il notamment contrôler ?
\end{questions}
\end{exercice}
\end{minipage}
\par\noindent\begin{minipage}{\linewidth}
\begin{exercice}[3]{L’air humide}
\begin{questions}
\item Pourquoi la composition de l’air sec ne décrit-elle pas exactement toute atmosphère humide ?
\item Peut-on additionner 78 \%, 21 \% et une forte part de vapeur d’eau sans réajuster ?
\end{questions}
\end{exercice}
\end{minipage}
\par\noindent\begin{minipage}{\linewidth}
\begin{exercice}[3]{Chercher une fuite}
\begin{questions}
\item Le piston reste fixe, mais la pression d’un gaz comprimé baisse lentement. Proposer deux explications.
\item Comment les départager ?
\end{questions}
\end{exercice}
\end{minipage}
\par\noindent\begin{minipage}{\linewidth}
\begin{exercice}[3]{Doubler l’échantillon}
\begin{questions}
\item On dispose de $250\unite{mL}$ puis de $500\unite{mL}$ d’air sec dans les mêmes conditions. Calculer le dioxygène correspondant.
\item La proportion a-t-elle doublé ?
\end{questions}
\end{exercice}
\end{minipage}
\par\noindent\begin{minipage}{\linewidth}
\begin{exercice}[3]{Un modèle incomplet}
\pcFigureEx{figures/ex-air-incorrect.pdf}{Modèle proposé par l’élève, à critiquer.}
\begin{questions}
\item Un dessin représente l’air par des points identiques immobiles, tous au fond d’une boîte. Identifier deux limites.
\item Proposer un dessin mieux légendé.
\end{questions}
\end{exercice}
\end{minipage}
\pcPage{Défi du chapitre}{L’air, les gaz et la pression}
\phantomsection\label{defi-ex-2}
\pcDefiTitre{2}{Dimensionner un modèle}
\begin{questions}
\item On veut représenter $1\unite L$ d’air sec par 200 jetons. Combien de jetons de chaque catégorie prévoir ?
\item Quel volume réel chaque jeton code-t-il dans ce modèle ?
\item Donner deux limites de la représentation.
\end{questions}
\pcGrille
